% Prospective task/data protocol; no new experiment or admitted model is reported.
\section{Follow the action, then test the prediction}

Watch a ball pass from one hand to the other. Where were the hands? When did one
let go? When did the other catch it? A model might follow the movement and still
get the handover wrong. That is why the next proposed study needs more than a
picture that looks convincing. It needs observations that can disagree with the
prediction.

We aim to connect a stated intention with an action as closely as possible in a
defined task. Understanding the process between them should make that connection
more specific. The test must still say whose movement it concerns, under which
conditions, how much error is tolerated and which independent observation checks
the result. It cannot turn success in one task into a complete map of the mind.

\subsection{Three recording views of the same event}

The scalp recording, arm sleeves and left/right tripod views introduced in
\emph{From a reach to a testable prediction} would have to agree on which event
they observed. The two surface-EMG sleeves would record one arm each; scalp EEG
would record scalp electrical activity, EMG muscle electrical activity, and
calibrated cameras the hands and arms. Their sites, units, capacity, clocks and
delays need checking before a handover can be called a shared observation.
OpenBCI or an equivalent configurable system remains an unacquired, unqualified
candidate.

For this hand task, the motion-capture relay-composite would join the left and
right observations. Earlier top and bottom dot maps describe another proposed
configuration. Each needs its own calibration and visibility checks. An unseen
finger remains unseen even if a model can draw a plausible replacement; the
display must keep that prediction separate from measured points.

The virtual body would then predict a response for comparison with the measured
movement. Its joints, muscles and contacts need appropriate reference evidence.
EMG does not directly measure force or incoming sensory-nerve activity. If the
claim concerns contact or force, the experiment needs an observation that can
check that claim.

\subsection{Let the task ask the next question}

Start with rapid hand rubbing. Contact and movement happen together, so matched
conditions are needed to find what each contributes. Record the requested
knuckle movement and the stretch with interlaced fingers, palms turned outward
and arms extended as distinct actions. The cold-hands example explains the
rubbing; it does not establish a cold-exposure experiment or a measured sensation.

Then pass a juggling ball between the hands---the proposed hot-potato task---and
progress to the preferred two-ball juggling task. One hand's release changes what
the other must do. The model now has to account for both limbs, the object and
the timing of a catch. Check it against simpler predictions from the cue or
previous movement; those may explain an apparent success without useful brain
information.

The boxing stage would distinguish stated punch types from independently
observed movement. It would become NeuroBoxer, a proposed virtual game. A
prime-era Mike Tyson opponent gives the game a recognizable aim; it supplies no
physical opponent or guarantee about injury. Boxing also uses the hands and
wrists for stabilization. Footage of an expert can show movement and tactics,
subject to its rights and camera limits. It cannot supply that expert's brain
signals, muscle signals or force measurements.

Toy archery adds another relationship. One hand holds while the other draws and
releases; the fingers change as the arrow leaves. The proposed apparatus has
arrows with small, cup-shaped suction tips and an over-arm pouch with several
arrows. NeuroArcher would turn the measured
coordination into action on a virtual target. A hit in the game still needs the
appropriate test before it can mean anything about biological control.

The complete snap described in \emph{From a reach to a testable prediction} tests another
relationship: preparation, contact and return must all agree with the specified
movement. The cup reach remains the simpler precision comparison. Both need
organism-specific anatomy, contact, timing and real-reference evidence; a
human-looking animation cannot establish that a monkey can perform the snap.

\subsection{Hands first; a brain contribution needs its own test}

The first game milestone would use measured body motion. A later milestone would
ask whether permitted EEG observations can control a specified virtual action
without using withheld muscle or movement targets. Eye activity, small covert
muscle contractions, residual movement, cues and information from the future
could all make a prediction look better. Those alternatives need declared
controls. Looking hands-free is not enough.

Keep the original brain-only, body-only and paired comparisons, including the
simple prediction that the target stays as it was. Set the allowed history and
the original success rules before fitting anything. Keep required people,
sessions and tasks separate before making time windows, and learn preprocessing
from training data only. A pretrained model's earlier exposure must be checked
before a recording can serve as independent final confirmation.

Recorded body observations could also predict a stated recorded EEG feature.
That reverse prediction is conditional: several brain states might fit the same
movement. It does not deliver a sensation. Qualified data and simpler baselines
come before paired fitting or fine-tuning. Reinforcement learning would need
specified actions, resulting transitions and rewards, rather than merely a
prediction error. A reward inside a simulation is not a biological result.
Final held-out tests stay outside training and selection; a running model stays
frozen until a candidate update is independently evaluated.

\subsection{A recording device does not become a writer}

The output boundary in \emph{Write sensation. Enact action} applies to these
tasks too. EEG and surface EMG record; neuromuscular electrical stimulation
(NMES) and functional electrical stimulation (FES) can cause contractions or
functional movement. A body-output experiment or sensory writer would need its
own qualified interface and permission. No spinal-writing interface is qualified here,
and these physical additions remain outside the acquisition and virtual-work
budget.

A proposed blindfolded condition removes visual feedback; it does not prove
neural control or learning. To ask what was retained, remove trainer information
while keeping the necessary movement bridge and comparable task conditions.
Check what the trainee can still do against the original controls. These
recordings, writer experiments and retention tests have not been performed
here. The current LINK failures remain failures; the sequence gives the next
paired study questions it would have to answer.
